% Teil II: der durchgerechnete Weg.

% Gedankenschritt: Die Schleife erhöht N, bis die Differenz eingeholt ist;
% das Restglied ~1/N erklärt, warum N in der 1000er-Größenordnung liegt.

\subsection*{Gegeben und gesucht}
(a) Differenz $\tfrac{\pi^2}6-\sum_{k=1}^{100}\frac1{k^2}$ in Python
berechnen; (b) das kleinste $N$ mit $\tfrac{\pi^2}6-\sum_{k=1}^{N}\frac1{k^2}<0.001$ finden.

\subsection*{a) Differenz bis $N=100$}
\begin{verbatim}
import math
diff = math.pi**2/6 - sum([1/k**2 for k in range(1, 101)])
print(diff)
\end{verbatim}
Das ergibt n\"aherungsweise
\begin{equation}
  \frac{\pi^2}6-\sum_{k=1}^{100}\frac1{k^2}
  \approx 0.0099502 .
\end{equation}

\subsection*{b) kleinstes $N$ mit Differenz $<0.001$}
\begin{verbatim}
import math
for N in range(1, 100000):
    diff = math.pi**2/6 - sum([1/k**2 for k in range(1, N+1)])
    if diff < 0.001:
        print(N, diff)
        break
\end{verbatim}
Das kleinste solche $N$ ist
\begin{equation}
  N=1000,\qquad \text{Differenz }\approx0.00099950.
\end{equation}

\begin{ergebnis}
\begin{align*}
  \text{a)}&\ \frac{\pi^2}6-\sum_{k=1}^{100}\frac1{k^2}\approx0.0099502,\\
  \text{b)}&\ N=1000\ \text{liefert erstmals eine Differenz }<0.001.
\end{align*}
\end{ergebnis}

\begin{kontrolle}
\textbf{Warum ist $N\sim10^3$ plausibel? (Grenzfall/Gr\"o\ss{}enordnung)}
Die Reihe konvergiert langsam wie $1/k^2$; das Restglied ist grob
\[
  \sum_{k>N}\frac1{k^2}\approx\int_N^\infty\frac{\mathrm dx}{x^2}
  =\frac1N .
\]
Also ist der Schwellwert $0.001\approx1/N$ erst bei $N\sim1000$
erreicht --- konsistent mit dem gefundenen $N=1000$. (Bei $N=100$ st\"unde
noch $\approx0.01$ Rest, zehnmal zu viel.)

\textbf{Off-by-One im Code (der eigentliche Test):} Das gefundene
$N=1000$ liegt \emph{jenseits} von $100$; eine Schleife bis nur $999$
(wie ``range(1,1000)``, oberer Index exklusiv) würde den Wert $N=1000$
nie erreichen und nie ausbrechen. Deshalb muss die Obergrenze der
Schleife deutlich über dem erwarteten $N$ liegen --- hier bis $10^5$.
Das ist genau das ``range(endet nicht am oberen Wert``-Prinzip., der hier
genau passt.
\end{kontrolle}

\begin{merksatz}
Iterativ $N$ erh\"ohen und die Differenz $\tfrac{\pi^2}6-$Teilsumme mit der
Schwelle vergleichen; wegen des langsamen $1/N$-Abklingens des
Restglieds liegt die Schwelle $0.001$ erst in der $10^3$-Gr\"o\ss{}en-
ordnung. \texttt{range(1,N+1)} und \texttt{k**2} sind die richtigen
Bausteine.
\end{merksatz}
